Interpreting the findings operationally: All results reported below are strictly conditional on the mathematical assumptions and behavioral parameters synthesized within the simulation model (as detailed in Section 1.5). They represent the optimized behavior of the HuDRL architecture, not definitive, direct predictions of real-world ecological or psychological outcomes. The answers to the research questions are therefore fundamentally comparative and exploratory. Within these operational bounds, the empirical results generated by the simulation suggest specific answers to the three research questions posed at the outset of this study, heavily emphasizing the performance of the HuDRL system:

\paragraph{RQ1: Impact of Behavioral Parameters.}

The global sensitivity analysis and comparative scenario testing demonstrated within the model that variations in household behavioral parameters fundamentally dictated policy stability. Specifically, the simulation revealed that the computational Cost of Effort'' was the dominant algorithmic inhibitor of compliance, overriding high Attitude'' scores \citep{Zheng2022}. This mathematically explained why simulated passive IEC campaigns failed; increasing awareness (Attitude) did not compensate for the high logistical friction (PBC) in areas parameterized like Poblacion \citep{Moeini2023}. Consequently, simulated policies that did not forcefully reduce this cost—either through subsidizing convenient infrastructure or deploying high-intensity enforcement to make non-compliance mathematically costlier—consistently failed to sustain compliance, regardless of public awareness levels \citep{Chen2023, Badua2022}.

\paragraph{RQ2: Optimal Resource Allocation Ratio.}

The simulation determined that there was no single static golden ratio'' for resource allocation. Instead, the optimal allocation was \textit{dynamic and state-dependent}. The HuDRL agent demonstrated that the most mathematically efficient strategy involved a phased approach: an initial heavy concentration (up to 90\%) on Enforcement and Incentives to break modeled resistance, followed by a gradual pivot toward IEC and maintenance spending once the heuristic tipping point'' of 7\% compliance was reached \citep{Nishimura2022}. This challenged the static 20\%/30\%/50\% splits often found in annual appropriation ordinances, demonstrating that allocation ratios had to evolve quarter-by-quarter based on real-time simulated compliance states \citep{Abdallah2021, Torkayesh2021}.

\paragraph{RQ3: The Superior Dynamic Strategy.}

The Heuristic-Guided PPO algorithm emerged as the dynamic strategy yielding the highest simulated cost-benefit ratio. While the Status Quo'' achieved only 12.5\% global compliance and Pure Enforcement'' collapsed to near-zero in key areas, the DRL-driven ``Sequential Saturation'' strategy successfully triggered the heuristic graduation of all barangays within 12 quarters without exceeding the strict 1.5 million peso budget constraint \citep{Jimenez2025}. This confirmed that, under the model's assumptions, the highest compliance per peso was achieved not by spending more, but by spending sequentially—concentrating capital to trigger simulated social norms that subsequently sustained behavior for free \citep{ZhengAI2022}. For an exhaustive quarter-by-quarter breakdown of this strategy's financial allocations and resulting compliance across all barangays, refer to the tabulated results in Appendix \ref{sec:hudrl_results}.

The comparative analysis and subsequent mechanistic investigation of the simulation results yield insights into the potential optimization of solid waste management (SWM) policies in the Municipality of Bacolod. The primary structural finding is that the conventional ``Status Quo'' approach---characterized by an equitable distribution of resources across all barangays and a heavy emphasis on Information, Education, and Communication (IEC)---is mathematically inefficient and insufficient to achieve municipality-wide compliance within the simulation. While this strategy maintained order in smaller, culturally cohesive community proxies, it consistently failed to overcome the high behavioral friction programmed into densely populated urban center proxies. This simulated failure was not due to a deficiency in total funding, but rather a strategic error in resource dilution; spreading the budget thinly prevented the Mayor Agent from reaching the critical enforcement intensity required to break the inertia of non-compliance in high-resistance zones.

Furthermore, the simulation demonstrates that single-instrument policies, whether purely punitive or purely incentive-based, are equally flawed when applied uniformly across a heterogeneous environment. Simulated pure enforcement eventually collapsed due to its inability to sustain positive community attitudes, while pure incentives failed due to the diminishing per-capita value of rewards in large populations. In contrast, the DRL agent’s Sequential Saturation'' strategy provides a robust computational solution to this dilemma. By abandoning the heuristic of simultaneous equality in favor of sequential efficacy, the agent successfully leveraged the modeled Graduation Effect.'' This mechanism utilized simulated social norms ($w_{sn} \approx 0.35$) to sustain compliance in stabilized areas while mathematically concentrating capital to overcome the structural inertia of the most resistant urban targets.

The evaluation of the Agent-Based Model through Global Sensitivity Analysis (GSA) provides a robust computational foundation for these optimized policies. The primary discovery of this evaluation is the overwhelming dominance of the Convenience Barrier''; with the Cost of Effort ($c_{\text{effort}}$) accounting for 79\% of output variance, the study computationally confirms that logistical friction is the single greatest structural inhibitor of waste segregation \citep{Yazawa2025}. Modeled policies that do not actively reduce this friction are statistically likely to fail regardless of simulated public support. This is further illustrated by the near-zero sensitivity of the Attitude ($w_a$) parameter, quantifying a distinct computational Value-Action Gap'' where awareness alone is insufficient to bridge the divide between intent and action \citep{Zhao2022}.

Ultimately, the HuDRL agent’s preference for localized enforcement and incentives represents a sophisticated mathematical response to the socio-technical barriers in Philippine waste management. It is important to reiterate the precise bounds of these findings: given the computational assumptions and static budget limits of the simulation, the model suggests that an adaptive, phased allocation of funds (Sequential Saturation) is the most optimal operational strategy. It cannot definitively forecast full real-world ecological shifts or human psychological changes. However, as an operational optimization tool, the RL+ABM framework provides a rigorously tested, data-driven roadmap toward a resilient and self-sustaining compliance ecosystem, demonstrating its immense utility as a decision-support mechanism for local government units.